% ============================================================
% 片段：s02-sca  对应原文 PDF 第 5--12 页
% 内容：§2 超共形代数与手征环（2.1--2.5），含编号公式 (2.1)--(2.42)
% 译者：t-s02   审校状态：初稿
% ============================================================

% ==== 原文 PDF 第 5 页 ====
\section{超共形代数与手征环}
\label{sec:2}

\begin{keytakeaways}
  \begin{itemize}[leftmargin=1.4em,itemsep=2pt,topsep=2pt]
    \item \textbf{物理图像}：镜像对称的代数起点是 $\mathcal{N}=(2,2)$ 超共形代数；它由左手与右手各一份 $\mathcal{N}=2$ 超共形代数组成，带横线的流属于右手（反全纯）扇区。
    \item \textbf{数学工具}：拓扑弦的 A 模型与 B 模型并不是完整的超共形场论（SCFT），而是把 $\mathcal{N}=(2,2)$ 表示的态空间\emph{截断到手征环}后得到的理论 \cite{ref5}；因此手征环的乘法与三点函数才是拓扑弦的物理数据。
    \item \textbf{关键结论}：手征主态满足 $h=q/2$（BPS 界），手征环在 OPE 下封闭且有限（$h\le c/6$）；由荷 $(1,1)$ 算符张成的形变子环上带有平坦的 Gauss--Manin 联络，其几何实现就是 B 模型一侧霍奇结构变化中的联络。
  \end{itemize}
\end{keytakeaways}

镜像对称（mirror symmetry）起源于 $\mathcal{N}=(2,2)$ 超共形代数（superconformal algebra）的表示；所谓 $\mathcal{N}=(2,2)$ 指的是两份 $\mathcal{N}=2$ 超共形代数，即左手移动与右手移动各一份，其中单份 $\mathcal{N}=2$ 代数将在下文讨论。由镜像对称互相等同的拓扑弦 A 模型与 B 模型，对应于把 $\mathcal{N}=(2,2)$ 超共形代数表示的态截断到手征环（chiral ring）\cite{ref5}。以下内容基于文献 \cite{ref27,ref29,ref6,ref30,ref31,ref10,ref32}。

% [Tutor Note: 物理图像] A 模型与 B 模型只是 $\mathcal{N}=(2,2)$ 表示的手征环截断
\begin{yizhu}[物理图像]
原文特意区分了两件事：$\mathcal{N}=(2,2)$ SCFT 的完整态空间，以及它在拓扑弦中实际用到的\emph{手征环截断}。只有截断后的环乘法 (2.18) 与三点函数 (2.22) 才是拓扑弦的物理数据；更高亏格的信息来自把该截断理论与引力耦合（§3），而不是来自完整 SCFT 的其余态。阅读时不要把 “$\mathcal{N}=(2,2)$ 的表示论” 与 “手征环” 等量齐观：前者是代数框架，后者才是拓扑弦的态空间。
\end{yizhu}

% [Tutor Note: 易错点] 单份 $\mathcal{N}=2$ 与 $\mathcal{N}=(2,2)$ 的记号切换点在第 8 页
\begin{yizhu}[易错点]
(2.1)--(2.17) 全部只涉及单份代数 $(T,G^{\pm},J)$；直到第 8 页说明“后者的流记作 $\bar T,\bar G^{\pm},\bar J$”时才引入第二份代数。因此 (2.14) 中的 $J_{0}$ 是单份 $U(1)$ 流的零模，而 (2.22) 中的 “荷 $(1,1)$” 是左右两个 $U(1)$ 荷组成的一对，二者不可混用；(2.17) 的 $c$ 是单份代数的中心荷，在 $\mathcal{N}=(2,2)$ 情形左右两份各有一个中心荷，不能把 $c/6$ 直接当作总中心荷给出的界。
\end{yizhu}

% [Tutor Note: 含糊点] “截断到手征环”的合法性由 §2.2 的闭性与有限性两条性质补足
\begin{yizhu}[含糊点]
原文只说态空间“截断到手征环”，没有解释为何可以截断。依据在 §2.2：手征主场的 OPE (2.16) 中指数满足 $h_{\psi_{i}}-h_{\phi}-h_{\chi}\ge 0$，故展开无奇性项；而 (2.17) 给出 $h\le c/6$，保证环有限。两者合起来才说明截断后的态空间在 OPE 下封闭、拓扑弦的关联函数良定义。可把 §2.2 当作对本段这句话的补证来读。
\end{yizhu}

\subsection{\texorpdfstring{$\mathcal{N}=2$}{N=2} 超共形代数}
\label{sec:2.1}

\begin{keytakeaways}
  \begin{itemize}[leftmargin=1.4em,itemsep=2pt,topsep=2pt]
    \item \textbf{物理图像}：$\mathcal{N}=2$ 超共形代数由能动量张量 $T(z)$（共形权 $2$）、两条超流 $G^{\pm}(z)$（共形权 $3/2$，$U(1)$ 荷为 $\pm1$）以及 $U(1)$ 流 $J(z)$ 生成。
    \item \textbf{数学工具}：全部信息可写成流的 OPE (2.1)--(2.6)，或等价地写成模式代数 (2.9)；边界条件 (2.7) 用一个连续参数 $a$ 统一描述，$a$ 的整数与半整数取值分别给出反周期与周期条件。
    \item \textbf{关键结论}：不同 $a$ 给出的代数彼此同构；该同构在态空间上诱导谱流，把 Neveu--Schwarz 扇区与 Ramond 扇区的态连续联系起来，从而诱导时空超对称。
  \end{itemize}
\end{keytakeaways}

$\mathcal{N}=2$ 超共形代数由能动量张量 $T(z)$（共形权 $h=2$）、两条反对易的超流 $G^{\pm}(z)$（共形权 $3/2$）以及一个 $U(1)$ 流 $J(z)$ 生成；$G^{\pm}(z)$ 在该 $U(1)$ 流下的荷为 $\pm$，而 $z$ 是二维世界面上的局部坐标。这些流满足如下算符乘积展开（operator product expansion，OPE）：
\begin{equation}\tag{2.1}
G^{\pm}(z)G^{\mp}(w)=\frac{\frac{2}{3}c}{(z-w)^{3}}\pm\frac{2J(w)}{(z-w)^{2}}+\frac{2T(w)\pm\partial_{w}J(w)}{(z-w)}+\ldots,
\end{equation}
\begin{equation}\tag{2.2}
J(z)G^{\pm}(w)=\pm\frac{G^{\pm}(w)}{(z-w)}+\ldots,
\end{equation}

% ==== 原文 PDF 第 6 页 ====
\begin{equation}\tag{2.3}
J(z)J(w)=\frac{\frac{1}{3}c}{(z-w)^{2}}+\ldots,
\end{equation}
\begin{equation}\tag{2.4}
T(z)J(w)=\frac{J(w)}{(z-w)^{2}}+\frac{\partial_{w}J(w)}{(z-w)}+\ldots,
\end{equation}
\begin{equation}\tag{2.5}
T(z)G^{\pm}(w)=\frac{\frac{3}{2}G^{\pm}(w)}{(z-w)^{2}}+\frac{\partial_{w}G^{\pm}(w)}{(z-w)}+\ldots,
\end{equation}
\begin{equation}\tag{2.6}
T(z)T(w)=\frac{\frac{1}{2}c}{(z-w)^{4}}+\frac{2T(w)}{(z-w)^{2}}+\frac{\partial_{w}T(w)}{(z-w)}+\ldots,
\end{equation}
其中省略号表示在 $z\to w$ 极限下添加的正规项，$c$ 是中心荷（central charge）。

% [Tutor Note: 易错点] (2.1) 的 $\pm$ 号与 (2.7) 的 $\mp$ 号属于同一套符号约定
\begin{yizhu}[易错点]
(2.1) 中 $G^{\pm}G^{\mp}$ 的 $\pm$ 同时作用于 $2J(w)/(z-w)^{2}$ 与 $\partial_{w}J(w)/(z-w)$ 两项：取上面符号时这两项都带正号，取下面符号时同时变号（首项 $\frac{2}{3}c/(z-w)^{3}$ 与符号无关）。类似地 (2.7) 中 $G^{+}$ 取 $\mathrm{e}^{-2\pi ia}$、$G^{-}$ 取 $\mathrm{e}^{+2\pi ia}$。这三处若各自独立约定，就会与 (2.9) 的 $[J_{n},G^{\pm}_{m\pm a}]=\pm G^{\pm}_{n+m\pm a}$ 矛盾；核对时可先看 (2.3) 的 $\frac{1}{3}c$ 与 (2.6) 的 $\frac{1}{2}c$ 是否与 (2.9) 首两行的 $c/3$、$c/12$ 对应。
\end{yizhu}

必须对超流 $G^{\pm}(z)$ 施加的边界条件可概括为
\begin{equation}\tag{2.7}
G^{\pm}(\mathrm{e}^{2\pi i}z)=-\mathrm{e}^{\mp 2\pi ia}G^{\pm}(z),
\end{equation}
其中连续参数 $a$ 的取值范围为 $0\le a<1$。对于整数与半整数 $a$，分别得到反周期与周期边界条件，它们对应于 Ramond 与 Neveu--Schwarz 扇区。
% [原文疑误: 扇区配对印反] 由 (2.8) 的模展开，$a\in\mathbb{Z}$ 时 $G^{\pm}$ 的模指标为半整数（反周期），对应 Neveu--Schwarz 扇区；$a\in\mathbb{Z}+\frac{1}{2}$ 时模指标为整数（周期），对应 Ramond 扇区。原文 “Ramond and Neveu--Schwarz, respectively” 与 (2.7)--(2.8) 的约定相反，照录并加此注。

% [Tutor Note: 原文疑误] (2.7) 之后一句把 Ramond 与 Neveu--Schwarz 的配对写反了
\begin{yizhu}[原文疑误]
按 (2.7) 直接验证：$a=0$ 时 $G^{\pm}(\mathrm{e}^{2\pi i}z)=-G^{\pm}(z)$ 为反周期，对应 Neveu--Schwarz 扇区；$a=\frac{1}{2}$ 时 $-\mathrm{e}^{\mp i\pi}=+1$ 为周期，对应 Ramond 扇区。这与 (2.8) 中 $a$ 进入模指标 $n\pm a$ 的方式一致（$a=0$ 给出半整数模即反周期）。因此原文“反周期、周期”与“Ramond、Neveu--Schwarz”的按序配对应当互换；译文照录原文，请以 (2.7)--(2.8) 的计算为准。
\end{yizhu}

这些流可按傅里叶模式展开
\begin{equation}\tag{2.8}
T(z)=\sum_{n}\frac{L_{n}}{z^{n+2}},\qquad
G^{\pm}(z)=\sum_{n}\frac{G^{\pm}_{n\pm a}}{z^{n\pm a+\frac{3}{2}}},\qquad
J(z)=\sum_{n}\frac{J_{n}}{z^{n+1}}.
\end{equation}

$\mathcal{N}=2$ 超共形代数既可用流的算符乘积展开表示，也可用其模式的交换关系表示。后者为
\begin{equation}\tag{2.9}
\begin{aligned}
[L_{m},L_{n}]&=(m-n)L_{m+n}+\frac{c}{12}m(m^{2}-1)\delta_{m+n,0},\\
[J_{m},J_{n}]&=\frac{c}{3}m\delta_{m+n,0},\\
[L_{n},J_{m}]&=-mJ_{m+n},\\
[L_{n},G^{\pm}_{m\pm a}]&=\Big(\frac{n}{2}-(m\pm a)\Big)G^{\pm}_{m+n\pm a},\\
[J_{n},G^{\pm}_{m\pm a}]&=\pm G^{\pm}_{n+m\pm a},\\
\{G^{+}_{n+a},G^{-}_{m-a}\}&=2L_{m+n}+(n-m+2a)J_{n+m}+\frac{c}{3}\Big((n+a)^{2}-\frac{1}{4}\Big)\delta_{m+n,0}.
\end{aligned}
\end{equation}

% [Tutor Note: 易错点] (2.9) 的 $c/12$、$c/3$ 与最后一条的指标可用 (2.14) 自检
\begin{yizhu}[易错点]
(2.9) 中 $[L_{m},L_{n}]$ 的 $c/12$、$[J_{m},J_{n}]$ 的 $c/3$ 分别来自 (2.6) 的 $c/2$ 与 (2.3) 的 $c/3$；最后一条的位移项 $\frac{c}{3}\big((n+a)^{2}-\frac{1}{4}\big)$ 则来自 (2.1) 的 $\frac{2}{3}c$。一个常用自检：取 $a=\frac{1}{2}$、$n=-1$、$m=1$，最后一条给出
$\{G^{+}_{-1/2},G^{-}_{1/2}\}=2L_{0}+(-1-1+1)J_{0}=2L_{0}-J_{0}$，
正是 (2.14) 中被积式。若把 $(n-m+2a)$ 误记成 $(m-n+2a)$，这一自检立即失败。
\end{yizhu}

对连续参数 $a$ 的每个取值所得的代数彼此同构。这一同构在态上诱导出一个操作，称为\emph{谱流}（spectral flow）。该操作表明 Ramond（R）扇区与 Neveu--Schwarz（NS）扇区等价：两个扇区中的态可通过谱流连续地相互联系。此外，由于 NS 与 R 扇区分别给出时空玻色子与时空费米子，代数的这一同构诱导出时空超对称。

% [Tutor Note: 物理图像] 谱流是代数同构在态空间上的实现，原文略去了具体机制
\begin{yizhu}[物理图像]
原文只声明“各 $a$ 的代数同构”并称之为谱流，略去了同构的实现方式：该同构由代数的一个分次自同构给出，它把 $G^{\pm}$ 的模式指标整体平移 $\theta$（即 $G^{\pm}_{r}\mapsto G^{\pm}_{r\pm\theta}$），同时把 $L_{0}$ 与 $J_{0}$ 的本征值作相应平移，因此同一代数内 NS 扇区的态可连续形变到 R 扇区的态。由此，NS 与 R 分别对应时空玻色子与费米子这一事实，才转化为“该同构诱导时空超对称”的结论。原文对这一步只给出一句话，这正是 §2.1 中最需要自行补算的地方。
\end{yizhu}

% ==== 原文 PDF 第 7 页 ====
\subsection{手征环}
\label{sec:2.2}

\begin{keytakeaways}
  \begin{itemize}[leftmargin=1.4em,itemsep=2pt,topsep=2pt]
    \item \textbf{物理图像}：手征主态满足 $h=q/2$，是 BPS 界的二维类比；它们在 OPE 下给出的手征环无奇点且有限，是拓扑弦可截断的态空间。
    \item \textbf{数学工具}：用 $\{G^{-}_{1/2},G^{+}_{-1/2}\}$ 与 $\{G^{-}_{3/2},G^{+}_{-3/2}\}$ 的半正定性分别推出 $h\ge q/2$ 与 $h\le c/6$。
    \item \textbf{关键结论}：$\mathcal{N}=(2,2)$ 情形共有 $(c,c)$、$(a,c)$、$(a,a)$、$(c,a)$ 四个有限环，后两者是前两者的电荷共轭；环结构由 (2.18) 的结构常数给出。
  \end{itemize}
\end{keytakeaways}

$\mathcal{N}=2$ 超共形代数的表示论带有一种有趣的附加结构：产生最高权态的算符中存在一个有限子扇区，它额外承载一个环结构，下面加以讨论。

代数的幺正不可约表示可由最高权态出发、用产生算符作用而建立；这些产生算符即 (2.8) 中负指标的模式。类似地，所有正指标的模式可视为湮灭算符，它们降低态的 $L_{0}$ 本征值。于是最高权态满足
\begin{equation}\tag{2.10}
L_{n}|\phi\rangle=0,\qquad G^{\pm}_{r}|\phi\rangle=0,\qquad J_{m}|\phi\rangle=0,\qquad n,r,m>0.
\end{equation}
零指标模式 $L_{0}$ 与 $J_{0}$ 可用来以其本征值标记态：
\begin{equation}\tag{2.11}
L_{0}|\phi\rangle=h_{\phi}|\phi\rangle,\qquad J_{0}|\phi\rangle=q_{\phi}|\phi\rangle.
\end{equation}
在 Ramond 扇区中还有模式 $G^{\pm}_{0}$。若一个态被这些模式湮灭，则称之为 Ramond 基态。最高权态由主场（primary field）$\phi$ 产生：
\begin{equation}\tag{2.12}
\phi|0\rangle=|\phi\rangle.
\end{equation}
我们关心的主场子集由\emph{手征主}（chiral primary）场构成。由它们产生的态还满足
\begin{equation}\tag{2.13}
G^{+}_{-1/2}|\phi\rangle=0.
\end{equation}
被 $G^{-}_{-1/2}$ 湮灭的主场将称为反手征主（anti-chiral primary）场。与反全纯流 $\bar G^{\pm}$ 的表示相结合，便引出 $(c,c)$、$(a,c)$、$(a,a)$ 与 $(c,a)$ 主场的概念；其中 $c$ 与 $a$ 分别代表手征与反手征，而这一对记号表示全纯与反全纯两个扇区中各自所加的条件。考察
\begin{equation}\tag{2.14}
\langle\phi|\{G^{-}_{1/2},G^{+}_{-1/2}\}|\phi\rangle=\|G^{+}_{-1/2}|\phi\rangle\|^{2}=\langle\phi|2L_{0}-J_{0}|\phi\rangle\ge 0,
\end{equation}
可得
\begin{equation}\tag{2.15}
h_{\phi}\ge\frac{q_{\phi}}{2},
\end{equation}
等号对手征态成立。手征主态的这一性质是物理态 BPS 界（BPS bound）的类比。现在考察两个手征主场 $\phi$ 与 $\chi$ 的算符乘积展开
\begin{equation}\tag{2.16}
\phi(z)\chi(w)=\sum_{i}(z-w)^{h_{\psi_{i}}-h_{\phi}-h_{\chi}}\psi_{i},
\end{equation}
$U(1)$ 荷相加，$q_{\psi_{i}}=q_{\phi}+q_{\chi}$，从而 $h_{\psi_{i}}\ge h_{\phi}+h_{\chi}$。因此算符乘积展开没有奇性项；当 $z\to w$ 时展开中幸存的项
% ==== 原文 PDF 第 8 页 ====
只有那些 $\psi_{i}$ 本身为手征主场的项。于是可见，手征主场在算符乘积展开下给出一个封闭的无奇点环。进一步，考察
\begin{equation}\tag{2.17}
\langle\phi|\{G^{-}_{3/2},G^{+}_{-3/2}\}|\phi\rangle=\langle\phi|2L_{0}-3J_{0}+\frac{2}{3}c|\phi\rangle\ge 0,
\end{equation}
可以证明该环的有限性，可见手征主场的共形权被 $c/6$ 界住。

% [Tutor Note: 易错点] (2.14) 的 $2L_{0}-J_{0}$ 完全由 (2.9) 最后一条的指标约定决定
\begin{yizhu}[易错点]
(2.14) 的算符组合应直接用 (2.9) 核对：取 $a=\frac{1}{2}$、$n=-1$、$m=1$，得
$\{G^{+}_{-1/2},G^{-}_{1/2}\}=2L_{0}+(n-m+2a)J_{0}=2L_{0}-J_{0}$。
若把 (2.9) 末行的 $(n-m+2a)$ 写成 $(m-n+2a)$，同一自检会得到 $2L_{0}+3J_{0}$，进而把 (2.15) 的界写成 $h\ge -3q/2$。同理 (2.17) 取 $a=\frac{1}{2}$、$n=-2$、$m=2$ 时，常数项为 $\frac{c}{3}\big(\frac{9}{4}-\frac{1}{4}\big)=\frac{2}{3}c$，这就是 (2.17) 中 $\frac{2}{3}c$ 的来源。
\end{yizhu}

% [Tutor Note: 含糊点] 从 (2.17) 到 $h\le c/6$ 之间省略了两步代数
\begin{yizhu}[含糊点]
原文说“考察 (2.17) 便知共形权被 $c/6$ 界住”，省去了两步：由半正定性得 $2h-3q+\frac{2}{3}c\ge 0$；再用手征条件 (2.15) 取等号即 $q=2h$，代入得 $2h-6h+\frac{2}{3}c\ge 0$，即 $h\le c/6$。注意这里用的是同一个手征态同时满足两条界：$h\ge q/2$ 给的是下界，$h\le c/6$ 给的是上界，二者结合才使手征环有限（每个荷级只有有限多个态）。
\end{yizhu}

% [Tutor Note: 易错点] 四个环与 $C^{k}_{ij}$ 的上指标要经 $\eta$ 降低后才等于 §2.4 的 $C_{ijk}$
\begin{yizhu}[易错点]
(2.18) 的 $\phi_{i}\phi_{j}=C^{k}_{ij}\phi_{k}$ 只是 $(c,c)$ 环的乘法；$(a,a)$ 与 $(c,a)$ 环的乘法由反全纯侧的条件定义，是前两者的电荷共轭。反手征主场的荷—权关系为 $h_{\psi}=-q_{\psi}/2$，与 (2.15) 的手征情形 $h_{\phi}=q_{\phi}/2$ 差一个负号，切勿混用。另外 $C^{k}_{ij}$ 的上指标 $k$ 需用拓扑度规降低后才得到 §2.4 中的 $C_{ijk}$（见 (2.23)），这一步的指标位置是后文矩阵表示 (2.24) 正确与否的关键。
\end{yizhu}

对于某些弦紧化的世界面描述，需要两份 $\mathcal{N}=2$ 超共形代数；它们通常称为左手与右手移动的，或全纯与反全纯的；后者的流将记作 $\bar T$、$\bar G^{\pm}$、$\bar J$。此时该代数称为 $\mathcal{N}=(2,2)$ 超共形代数。

在 $\mathcal{N}=(2,2)$ 超共形代数中现在有四个有限环，取决于左右移动代数中手征与反手征的不同组合；由此得到 $(c,c)$、$(a,c)$、$(a,a)$ 与 $(c,a)$ 环。可以看出后两者是前两者的电荷共轭。反手征主场的荷与共形权之间的关系变为 $h_{\psi}=-\frac{q_{\psi}}{2}$。以 $\phi_{i}$ 记手征主场的集合，指标 $i$ 遍历所有手征主场，则环结构可表述为如下形式\footnote{算符乘积的公式应理解为在关联函数内成立。}：
\begin{equation}\tag{2.18}
\phi_{i}\phi_{j}=C^{k}_{ij}\phi_{k}.
\end{equation}

\subsection{形变族}
\label{sec:2.3}

\begin{keytakeaways}
  \begin{itemize}[leftmargin=1.4em,itemsep=2pt,topsep=2pt]
    \item \textbf{物理图像}：共形场论的形变通过向作用量添加边缘算符实现；保持 $h=\bar h=1$ 的精确边缘算符对应 SCFT 模空间 $\mathcal{M}$ 上的坐标。
    \item \textbf{数学工具}：从荷 $(1,1)$、$h=\bar h=\frac{1}{2}$ 的手征主算符出发，先取 $G^{-}$ 的留数、再取 $\bar G^{-}$ 的留数，两步得到荷为零的 $(1,1)$ 算符 $\phi^{(2)}$。
    \item \textbf{关键结论}：形变作用量由 (2.21) 给出，形变空间即 SCFT 的模空间；上标记号借自拓扑场论的降格方程。
  \end{itemize}
\end{keytakeaways}

镜像对称是这样一个对称性：它把 $(a,c)$ 手征环的形变族（deformation family）与 $(c,c)$ 手征环的形变族联系起来。为使这一陈述更精确，下面讨论 $(c,c)$ 手征环的形变族。共形场论的形变通过向原作用量添加边缘算符（marginal operator）来实现；这些算符的共形权为 $h+\bar h=2$。下面研究无自旋算符，即 $h=1$、$\bar h=1$ 的算符。在对理论作微扰后仍保持其 $h=\bar h=1$ 共形权的算符称为\emph{精确边缘}（truly marginal）算符。这类算符可由手征主算符分两步构造。例如在 $(c,c)$ 环中，从电荷 $q=\bar q=1$、$h=\bar h=1/2$ 的算符出发，第一步构造
\begin{equation}\tag{2.19}
\phi^{(1)}(w,\bar w)=[G^{-}(z),\phi(w,\bar w)]=\oint\mathrm{d}z\,G^{-}(z)\phi(w,\bar w),
\end{equation}
它现在具有 $h=1$、$q=0$。下一步构造
\begin{equation}\tag{2.20}
\phi^{(2)}(w,\bar w)=\{\bar G^{-}(z),\phi^{(1)}(w,\bar w)\}=\oint\mathrm{d}\bar z\,\bar G^{-}(\bar z)\phi^{(1)}(w,\bar w),
\end{equation}
% [原文疑误: (2.20) 括号内写作 $\bar G^{-}(z)$] 同一式中 contour 积分的被积函数为 $\bar G^{-}(\bar z)$、积分微元为 $\mathrm{d}\bar z$，而交换子一侧却印成 $\bar G^{-}(z)$；按反全纯流的自变量约定应为 $\bar G^{-}(\bar z)$。照录原文并加此注。

% [Tutor Note: 原文疑误] (2.20) 中 $\bar G^{-}$ 的自变量前后不一致
\begin{yizhu}[原文疑误]
(2.20) 的交换子写作 $\{\bar G^{-}(z),\phi^{(1)}(w,\bar w)\}$，但同一式右端的积分是 $\oint\mathrm{d}\bar z\,\bar G^{-}(\bar z)\phi^{(1)}(w,\bar w)$。反全纯流应以 $\bar z$ 为自变量，故括号中应为 $\bar G^{-}(\bar z)$。这一不一致不影响 $h,\bar h$ 的计数，但动手计算留数时若把 $\bar G^{-}(z)$ 当成全纯对象处理，就会得到错误的共形权。
\end{yizhu}

% ==== 原文 PDF 第 9 页 ====
它具有 $h=\bar h=1$ 且荷为零，因此是精确边缘算符，可用来微扰理论的作用量
\begin{equation}\tag{2.21}
\delta S=t^{i}\int\phi_{i}^{(2)}+\bar t^{\bar{\imath}}\int\phi_{\bar{\imath}}^{(2)},\qquad i=1,\ldots,n,
\end{equation}
其中先验地也包括来自 $(a,a)$ 算符的形变，而 $n=\dim\mathcal{H}^{(1,1)}$ 表示理论的希尔伯特空间中由荷为 $(1,1)$ 的算符所产生态张成的子空间的维数。对 $(a,c)$ 手征环可作类似构造。上标记号借自拓扑场论；在拓扑场论中，类似构造给出二形式后代（two-form descendants），可用于微扰拓扑理论，参见例如 \cite{ref33}。以此方式构造的形变张成一个形变空间 $\mathcal{M}$，即 SCFT 的模空间。

% [Tutor Note: 易错点] 交换子与反对易子、以及 $t^{i}$ 与 $\bar t^{\bar{\imath}}$ 的指标位置
\begin{yizhu}[易错点]
(2.19) 用的是与全纯超流的\emph{交换子} $[G^{-}(z),\phi]$，(2.20) 用的是与反全纯超流的\emph{反对易子} $\{\bar G^{-},\phi^{(1)}\}$；两者都是 contour 积分（留数），但前者取 $\mathrm{d}z$、后者取 $\mathrm{d}\bar z$。此外 (2.21) 中全纯形变系数 $t^{i}$ 与反全纯系数 $\bar t^{\bar{\imath}}$ 的指标都在上方（后者还带横线），与 (2.22) 中 $C_{ijk}$ 的全下指标形成对比；不要把 $\bar t^{\bar{\imath}}$ 写成 $\bar t_{i}$。还有一处原文自身的记号不一致：(2.21) 之后写 $n=\dim\mathcal{H}^{(1,1)}$，而 (2.22) 写 $\dim\mathcal{H}^{1,1}$，同一空间带不带括号混用，译文照录。
\end{yizhu}

% [Tutor Note: 含糊点] 上标 $(2)$ 的含义被原文一句“借自拓扑场论”带过
\begin{yizhu}[含糊点]
原文只说上标记号借自拓扑场论、并指向 \cite{ref33}，未解释 $(2)$ 是什么。在拓扑场论中，从零形式算符 $O^{(0)}$ 出发，降格方程 $\mathrm{d}O^{(n)}=\{Q,O^{(n+1)}\}$ 依次给出一形式与二形式后代；$\int_{\Sigma}O^{(2)}$ 正是可以加进作用量的形变量。因此 $\phi_{i}^{(2)}$ 的上标表示它是二形式后代，这也解释了为何 (2.21) 可以直接把它对世界面积分。
\end{yizhu}

% [Tutor Note: 物理图像] 两步构造把 $(1,1)$ 荷的手征主算符变成精确边缘算符
\begin{yizhu}[物理图像]
构造的起点是荷 $(1,1)$、$h=\bar h=\frac{1}{2}$ 的手征主算符。第一步用 $G^{-}$ 把 $h$ 从 $\frac{1}{2}$ 提到 $1$、同时把 $q$ 从 $1$ 降到 $0$；第二步用 $\bar G^{-}$ 把 $\bar h$ 从 $\frac{1}{2}$ 提到 $1$、把 $\bar q$ 从 $1$ 降到 $0$。结果 $\phi^{(2)}$ 的共形权为 $(1,1)$、两侧荷均为零，恰好满足无自旋边缘算符的条件，可按 (2.21) 加入作用量。这正是“精确边缘算符 = 手征环中荷 $(1,1)$ 的元素经两次超流留数提升”的含义。
\end{yizhu}

\subsection{形变子环}
\label{sec:2.4}

\begin{keytakeaways}
  \begin{itemize}[leftmargin=1.4em,itemsep=2pt,topsep=2pt]
    \item \textbf{物理图像}：态丛 $\mathcal{V}$ 的纤维是理论基态的一个子集；基态总数在形变下不变，但按 $U(1)$ 荷的分解随模空间光滑变化。
    \item \textbf{数学工具}：用拓扑度规 $\eta_{im}$ 升降指标，把手征环乘法写成矩阵 $C_{i}$；插入 $\phi_{i}$ 等于 $\partial/\partial t^{i}$，故 $\nabla=\partial_{i}-C_{i}$ 是丛上的联络。
    \item \textbf{关键结论}：该联络平坦（见 (2.29)），称为 Gauss--Manin 联络，其在 B 模型一侧与霍奇结构的变化中定义的 Gauss--Manin 联络一致。
  \end{itemize}
\end{keytakeaways}

为研究 SCFT 之间等价的几何实现，引入形变空间 $\mathcal{M}$ 上的全纯向量丛 $V_{\mathbb{C}}$。其纤维对应于理论基态的一个子集。其重要性在于：理论的基态在形变空间上不变 \cite{ref9}。然而，存在一种在模空间上光滑变化的丛分解方式。该丛收集由荷为 $(1,1)$ 的算符张成的手征环子环所产生的态。该子环的一组基记作 $(\phi_{0},\phi_{i},(\phi_{D})_{i},(\phi_{D})_{0})$，$i=1,\ldots,n$。$\phi_{0}$ 是荷为 $(0,0)$ 的单位算符，$(\phi_{D})_{i}$ 是荷为 $(2,2)$ 的场，$(\phi_{D})_{0}$ 是荷为 $(3,3)$ 的元素。

球面上的三点函数是 3 个荷为 $(1,1)$ 的手征场的关联函数\footnote{由于 $U(1)$ 流中的反常，拓扑理论中关联函数里所有算符的荷之和必须为 $(3,3)$（参见例如 \cite{ref30,ref10}）。}：
\begin{equation}\tag{2.22}
C_{ijk}=\langle\phi_{i}\phi_{j}\phi_{k}\rangle,\qquad i,j,k=1,\ldots\dim\mathcal{H}^{1,1}.
\end{equation}
利用手征环关系 (2.18)，可作如下操作：
\begin{equation}\tag{2.23}
C_{ijk}=\langle\phi_{i}\phi_{j}\phi_{k}\rangle=\langle\phi_{i}C^{m}_{jk}(\phi_{D})_{m}\rangle=C^{m}_{jk}\langle\phi_{i}(\phi_{D})_{m}\rangle=C^{m}_{jk}\eta_{im},
\end{equation}
由此定义拓扑度规 $\eta_{im}$，它给出荷 $(1,1)$ 与荷 $(2,2)$ 算符之间的配对。该度规可用于升降指标，并把手征环的结构常数与三点函数等同起来\footnote{关于拓扑度规的更多细节，参见综述 \cite{ref33}。}。

% [Tutor Note: 易错点] (2.23) 的指标顺序依赖三点函数的全对称性
\begin{yizhu}[易错点]
(2.23) 把 $C_{ijk}$ 写成 $C^{m}_{jk}\eta_{im}$，而 (2.28) 中 $\phi_{i}$ 作用在 $|e_{j}\rangle$ 上给出的系数也是 $C_{ijk}$，即 $\phi_{i}\phi_{j}=C_{ijk}\phi^{k}$。两者一致要求 $C_{ijk}=C_{jki}=C_{kij}$（三点函数全对称）以及 $\eta_{im}=\eta_{mi}$；若忽略指标顺序，展开 (2.23) 的等式链就会差一个置换。核对方法：由 (2.24) 之前的 $\phi^{i}=\eta^{ij}(\phi_{D})_{j}$ 的对偶形式 $\phi_{i}=\eta_{im}\phi^{m}$ 重算一遍。
\end{yizhu}

% [Tutor Note: 易错点] (2.24)/(2.28) 矩阵的非零元与 $\delta^{j}_{i}$ 的方向
\begin{yizhu}[易错点]
两个矩阵的布局完全相同：第一行给出 $\phi_{i}\phi_{0}=\phi_{i}$（元 $\delta^{k}_{i}$，作用在第二列的 $\phi_{k}$ 上），第二行给出 $\phi_{i}\phi_{j}=C_{ijk}\phi^{k}$（元 $C_{ijk}$，作用在第三列的 $\phi^{k}$/ $|e^{k}\rangle$ 上），第三行给出 $\phi_{i}\phi^{j}=\delta^{j}_{i}\phi^{0}$（元 $\delta^{j}_{i}$，位于第四列、作用在 $\phi^{0}$ 上），第四行全零（$\phi_{i}\phi^{0}=0$）。$\delta^{j}_{i}$ 的指标顺序不可写反：列向量各分量的指标依次是 $0,j,j,0$，写成 $\delta^{i}_{j}$ 就与列向量的指标位置不符。
\end{yizhu}

% ==== 原文 PDF 第 10 页 ====
利用 $\phi^{i}=\eta^{ij}(\phi_{D})_{j}$ 以及 $\langle(\phi_{D})_{0}\rangle=1$、$\eta_{00}=1$，环结构现可写成矩阵形式
\begin{equation}\tag{2.24}
\phi_{i}
\begin{pmatrix}
\phi_{0}\\
\phi_{j}\\
\phi^{j}\\
\phi^{0}
\end{pmatrix}
=
\begin{pmatrix}
0 & \delta^{k}_{i} & 0 & 0\\
0 & 0 & C_{ijk} & 0\\
0 & 0 & 0 & \delta^{j}_{i}\\
0 & 0 & 0 & 0
\end{pmatrix}
\begin{pmatrix}
\phi_{0}\\
\phi_{k}\\
\phi^{k}\\
\phi^{0}
\end{pmatrix}.
\end{equation}
由该子环产生的态组织成一个向量丛
\begin{equation}\tag{2.25}
\mathcal{V}=\mathcal{H}^{0,0}\oplus\mathcal{H}^{1,1}\oplus\mathcal{H}^{2,2}\oplus\mathcal{H}^{3,3},
\end{equation}
其中 $\mathcal{H}^{(i,i)}$ 表示由荷为 $(i,i)$ 的算符产生的态所张成的希尔伯特空间子空间。丛的分解由电荷分次给出。该分次在模空间 $\mathcal{M}$ 上的变化及其几何实现，将是镜像对称研究中的中心主题。可利用算符--态对应从手征环算符的基得到向量丛的基。记 $|e_{0}\rangle\in\mathcal{H}^{0,0}$ 为荷 $(0,0)$ 的唯一基态（差一个标度），则 $\mathcal{V}$ 的基可如下得到：
\begin{equation}\tag{2.26}
|e_{i}\rangle=\phi_{i}|e_{0}\rangle,\qquad |e^{j}\rangle=\phi^{j}|e_{0}\rangle,\qquad |e^{0}\rangle=\phi^{0}|e_{0}\rangle.
\end{equation}
该丛上的度规由拓扑度规给出
\begin{equation}\tag{2.27}
\langle e_{a}|e^{b}\rangle=\delta^{b}_{a}\qquad a,b=0,\ldots,n.
\end{equation}
手征环在此基下的表示写成
\begin{equation}\tag{2.28}
\phi_{i}
\begin{pmatrix}
|e_{0}\rangle\\
|e_{j}\rangle\\
|e^{j}\rangle\\
|e^{0}\rangle
\end{pmatrix}
=
\underbrace{
\begin{pmatrix}
0 & \delta^{k}_{i} & 0 & 0\\
0 & 0 & C_{ijk} & 0\\
0 & 0 & 0 & \delta^{j}_{i}\\
0 & 0 & 0 & 0
\end{pmatrix}}_{:=C_{i}}
\begin{pmatrix}
|e_{0}\rangle\\
|e^{k}\rangle\\
|e^{k}\rangle\\
|e^{0}\rangle
\end{pmatrix}.
\end{equation}
这一结构对理解镜像对称将是关键的。就模空间 $\mathcal{M}$ 上的坐标而言，插入手征场 $\phi_{i}$ 等价于在模空间中作一无穷小位移，因而可由导数 $\frac{\partial}{\partial t^{i}}$ 得到。记右端矩阵为 $C_{i}$，则整个方程可读作丛上的一个联络 $\nabla=\partial_{i}-C_{i}$，它是平坦的\footnote{更多细节参见 \cite{ref19,ref9,ref34}。}：
\begin{equation}\tag{2.29}
[\nabla_{i},\nabla_{j}]=0.
\end{equation}
该联络称为 Gauss--Manin 联络，因为它在 B 模型一侧的几何实现与霍奇结构的变化中定义的 Gauss--Manin 联络一致。

% [Tutor Note: 物理图像] 基态总数受保护不变，按荷的分解却随模空间变化
\begin{yizhu}[物理图像]
原文说“基态在形变空间上不变 \cite{ref9}”，同时又说存在随模空间光滑变化的丛分解。这两句话并不矛盾：基态的总维数受 Witten 指标保护、不随形变改变，但各个荷级的基态会在不同荷级之间重排。因此 (2.25) 的分解 $\mathcal{V}=\mathcal{H}^{0,0}\oplus\cdots\oplus\mathcal{H}^{3,3}$ 随 $\mathcal{M}$ 光滑变化，而 (2.28) 的矩阵 $C_{i}$ 正是描述这种变化的联络。理解这一点，才能明白为何“基态不变”仍能给出非平凡的几何。
\end{yizhu}

% [Tutor Note: 含糊点] $[\nabla_{i},\nabla_{j}]=0$ 的平坦性来源被原文略去
\begin{yizhu}[含糊点]
原文直接断言 $\nabla=\partial_{i}-C_{i}$ 平坦并给出 (2.29)，未展开证明。其来源是：$\partial_{i}$ 来自模空间坐标的偏导（彼此交换），$C_{i}$ 来自手征环乘法的结合性与交换性，二者合起来使 $[\nabla_{i},\nabla_{j}]=0$。在 B 模型一侧，这一联络的几何实现正是霍奇结构变化中定义的 Gauss--Manin 联络，故原文随即给出该名称；至于反全纯方向的联络与度规的关系，属于特殊几何的内容，本章不展开。
\end{yizhu}

% ==== 原文 PDF 第 11 页 ====
\subsection{SCFT 的几何实现}
\label{sec:2.5}

\begin{keytakeaways}
  \begin{itemize}[leftmargin=1.4em,itemsep=2pt,topsep=2pt]
    \item \textbf{物理图像}：$\mathcal{N}=(2,2)$ SCFT 的几何实现是非线性 sigma 模型；靶空间须为凯勒流形，而 Ricci 平坦给出的正是卡拉比--丘流形。
    \item \textbf{数学工具}：作用量 (2.30) 由度规、$B$ 场、费米子动能与曲率四部分构成；费米子按 (2.31) 取 $K^{1/2}$ 或 $\bar K^{1/2}$ 的截面。
    \item \textbf{关键结论}：超对称变换 (2.32)--(2.37) 与流 (2.38)--(2.41) 给出 (2.1) 的几何实现，超对称变换可用荷写成 (2.42)。
  \end{itemize}
\end{keytakeaways}

$\mathcal{N}=(2,2)$ SCFT 的几何实现由非线性 sigma 模型提供 \cite{ref35,ref6}，另见综述 \cite{ref9,ref31,ref10}。这里的焦点是 $(a,c)$ 与 $(c,c)$ 环的几何实现。所考虑的非线性 sigma 模型是生活在二维黎曼面 $\Sigma$ 上的玻色子 $\phi$ 与费米子 $\psi,\chi$ 的场论，它们由超对称相联系，其中玻色场被看作某个靶空间 $X$ 的坐标，即 $\phi:\Sigma\to X$。为使理论具有 $\mathcal{N}=(2,2)$ 超对称，由玻色场张成的靶空间必须是凯勒流形（Kähler manifold），这又允许把其切丛分解为 $TX=T^{1,0}X\oplus T^{0,1}X$。理论的作用量为（如文献 \cite{ref6,ref10} 所给出）\footnote{这里的记号与文献 \cite{ref6,ref10} 的选择不同，以便更适合讨论高亏格拓扑弦；那里采用文献 \cite{ref19} 的记号。上标 $\pm$ 指超共形代数流的 $U(1)$ 荷，而表达式上方的横线指右手移动扇区。对费米子而言，左手移动费米子记作 $\psi$，右手移动费米子记作 $\chi$。}
\begin{equation}\tag{2.30}
\begin{aligned}
S=\;&\frac{1}{4\pi}\int_{\Sigma}\mathrm{d}^{2}z\Big(g_{i\bar{\jmath}}\big(\partial_{z}\phi^{i}\partial_{\bar{z}}\phi^{\bar{\jmath}}+\partial_{\bar{z}}\phi^{i}\partial_{z}\phi^{\bar{\jmath}}\big)+B_{i\bar{\jmath}}\big(\partial_{z}\phi^{i}\partial_{\bar{z}}\phi^{\bar{\jmath}}-\partial_{\bar{z}}\phi^{i}\partial_{z}\phi^{\bar{\jmath}}\big)\\
&+ig_{i\bar{\jmath}}\psi^{i}D_{z}\psi^{\bar{\jmath}}+ig_{i\bar{\jmath}}\chi^{i}D_{\bar{z}}\chi^{\bar{\jmath}}+R_{i\bar{k}j\bar{l}}\psi^{i}\psi^{\bar{k}}\chi^{j}\chi^{\bar{l}}\Big).
\end{aligned}
\end{equation}
记 $K$ 为 $\Sigma$ 上的典范丛，则费米子为如下丛的截面：
\begin{equation}\tag{2.31}
\begin{aligned}
\psi^{i}&\in\Gamma(K^{1/2}\otimes\phi^{*}T^{1,0}X), & \psi^{\bar{\imath}}&\in\Gamma(K^{1/2}\otimes\phi^{*}T^{0,1}X),\\
\chi^{i}&\in\Gamma(\bar{K}^{1/2}\otimes\phi^{*}T^{1,0}X), & \chi^{\bar{\imath}}&\in\Gamma(\bar{K}^{1/2}\otimes\phi^{*}T^{0,1}X).
\end{aligned}
\end{equation}
作用量中的协变导数是对这些丛取的。费米子 $\psi$ 与 $\chi$ 分别对应左手与右手移动扇区。若 Ricci 张量为零，则作用量共形不变；与凯勒条件一起，这给出卡拉比--丘靶空间。超对称由如下场变换给出：
\begin{equation}\tag{2.32}
\delta\phi^{i}=i\varepsilon^{-}\psi^{i}+i\bar\varepsilon^{-}\chi^{i},
\end{equation}
\begin{equation}\tag{2.33}
\delta\phi^{\bar{\imath}}=i\varepsilon^{+}\psi^{\bar{\imath}}+i\bar\varepsilon^{+}\chi^{\bar{\imath}},
\end{equation}
\begin{equation}\tag{2.34}
\delta\psi^{i}=-\varepsilon^{+}\partial\phi^{i}-i\bar\varepsilon^{-}\chi^{j}\Gamma^{i}_{jk}\psi^{k},
\end{equation}
\begin{equation}\tag{2.35}
\delta\psi^{\bar{\imath}}=-\varepsilon^{-}\partial\phi^{i}-i\bar\varepsilon^{+}\chi^{\bar{\jmath}}\Gamma^{\bar{\imath}}_{\bar{\jmath}\bar{k}}\psi^{\bar{k}},
\end{equation}
% [原文疑误: (2.35) 首项缺少 $\bar\partial$ 与带横指标的 $\phi^{\bar{\imath}}$] 与 (2.34) 的 $-\varepsilon^{+}\partial\phi^{i}$ 对照，反全纯费米子的变换首项应为 $-\varepsilon^{-}\bar\partial\phi^{\bar{\imath}}$；(2.35) 印作 $-\varepsilon^{-}\partial\phi^{i}$，与 (2.37) 的 $\bar\partial\phi^{\bar{\imath}}$ 结构不一致。照录原文并加此注。
\begin{equation}\tag{2.36}
\delta\chi^{i}=-\bar\varepsilon^{+}\bar\partial\phi^{i}-i\varepsilon^{-}\psi^{j}\Gamma^{i}_{jk}\chi^{k},
\end{equation}
\begin{equation}\tag{2.37}
\delta\chi^{\bar{\imath}}=-\bar\varepsilon^{-}\bar\partial\phi^{\bar{\imath}}-i\varepsilon^{+}\psi^{\bar{\jmath}}\Gamma^{\bar{\imath}}_{\bar{\jmath}\bar{k}}\chi^{\bar{k}},
\end{equation}
其中 $\varepsilon^{-},\bar\varepsilon^{-},\varepsilon^{+},\bar\varepsilon^{+}$ 分别是 $K^{-1/2},\bar{K}^{-1/2},K^{-1/2},\bar{K}^{-1/2}$ 的截面。流 (2.1) 的几何实现由下式给出：
\begin{equation}\tag{2.38}
T=-g_{i\bar{\jmath}}\partial_{z}\phi^{i}\partial_{z}\phi^{\bar{\jmath}}+\frac{1}{2}g_{i\bar{\jmath}}\psi^{i}\partial_{z}\psi^{\bar{\jmath}}+\frac{1}{2}g_{i\bar{\jmath}}\psi^{\bar{\jmath}}\partial_{z}\psi^{i},
\end{equation}

% [Tutor Note: 易错点] (2.30) 的 $1/4\pi$、$g/B$ 组合与曲率项的缩并
\begin{yizhu}[易错点]
(2.30) 的整体因子是 $\frac{1}{4\pi}$，不是 $\frac{1}{2\pi}$；$g$ 项写作对称组合 $\partial_{z}\phi^{i}\partial_{\bar{z}}\phi^{\bar{\jmath}}+\partial_{\bar{z}}\phi^{i}\partial_{z}\phi^{\bar{\jmath}}$，$B$ 项则是同两项的反对称组合，两者共用同一个 $B_{i\bar{\jmath}}$。曲率项 $R_{i\bar{k}j\bar{l}}\psi^{i}\psi^{\bar{k}}\chi^{j}\chi^{\bar{l}}$ 中四个费米子指标与曲率张量四个指标逐一缩并，注意第二个指标是 $\bar{k}$、第四个是 $\bar{l}$，不要按 $R_{i\bar{\jmath}k\bar{l}}$ 的习惯顺序写。
\end{yizhu}

% [Tutor Note: 易错点] (2.31) 中 $\psi^{\bar{\imath}}$ 属于 $K^{1/2}$ 而不是 $\bar K^{1/2}$
\begin{yizhu}[易错点]
(2.31) 中区分手征性的只有 $K$ 与 $\bar K$：$\psi^{i}$ 与 $\psi^{\bar{\imath}}$ 都取 $K^{1/2}$ 的截面，$\chi^{i}$ 与 $\chi^{\bar{\imath}}$ 都取 $\bar K^{1/2}$ 的截面。因此 $\psi^{\bar{\imath}}\in\Gamma(K^{1/2}\otimes\phi^{*}T^{0,1}X)$，横指标只表示靶空间的 $(0,1)$ 方向，与左手/右手无关。这正是脚注 6 所强调的记号：$\psi$ 是左手费米子、$\chi$ 是右手费米子。若把 $\psi^{\bar{\imath}}$ 误归入 $\bar K^{1/2}$，就会与 (2.30) 的 $D_{z}\psi^{\bar{\jmath}}$ 以及 (2.38) 的 $\partial_{z}\psi^{\bar{\jmath}}$ 矛盾。
\end{yizhu}

% [Tutor Note: 原文疑误] (2.35) 首项与相邻各式的 $\partial/\bar\partial$ 结构不相容
\begin{yizhu}[原文疑误]
(2.35) 印作 $\delta\psi^{\bar{\imath}}=-\varepsilon^{-}\partial\phi^{i}-\ldots$。对照 (2.34) 的 $\delta\psi^{i}=-\varepsilon^{+}\partial\phi^{i}-\ldots$ 与 (2.37) 的 $\delta\chi^{\bar{\imath}}=-\bar\varepsilon^{-}\bar\partial\phi^{\bar{\imath}}-\ldots$，反全纯费米子 $\psi^{\bar{\imath}}$ 的变换首项应含反全纯导数与带横指标的场，即 $-\varepsilon^{-}\bar\partial\phi^{\bar{\imath}}$；否则 (2.32)--(2.37) 的超对称代数与 (2.30) 中 $\psi$ 只与 $\partial_{z}$ 配对的动能项不相容。译文按原文照录，请在推演超对称代数时以 $\bar\partial\phi^{\bar{\imath}}$ 为准。
\end{yizhu}

% ==== 原文 PDF 第 12 页 ====
\begin{equation}\tag{2.39}
G^{+}=\frac{1}{2}g_{i\bar{\jmath}}\psi^{i}\partial_{z}\phi^{\bar{\jmath}},
\end{equation}
\begin{equation}\tag{2.40}
G^{-}=\frac{1}{2}g_{i\bar{\jmath}}\psi^{\bar{\jmath}}\partial_{z}\phi^{i},
\end{equation}
\begin{equation}\tag{2.41}
J=\frac{1}{4}g_{i\bar{\jmath}}\psi^{i}\psi^{\bar{\jmath}},
\end{equation}
类似地有 $\bar T,\bar G^{+},\bar G^{-},\bar J$。以超对称流之荷表述的超对称变换为\footnote{$Q^{\pm}=\frac{1}{2\pi i}\oint G^{\pm}\mathrm{d}z$，\quad $\bar Q^{\pm}=\frac{1}{2\pi i}\oint\bar G^{\pm}\mathrm{d}\bar z$。}：
\begin{equation}\tag{2.42}
\delta=\varepsilon^{+}Q^{+}+\varepsilon^{-}Q^{-}+\bar\varepsilon^{+}\bar Q^{+}+\bar\varepsilon^{-}\bar Q^{-}.
\end{equation}

% [Tutor Note: 物理图像] 流 (2.38)--(2.41) 把抽象代数写成靶空间几何量
\begin{yizhu}[物理图像]
(2.38)--(2.41) 是 (2.1) 的具体化：$T$ 由玻色动能与费米子动能组成，$G^{\pm}$ 形如 $\psi\partial\phi$（左右两种配对），而 $J=\frac{1}{4}g_{i\bar{\jmath}}\psi^{i}\psi^{\bar{\jmath}}$ 是费米子数流。三式的系数 $\frac{1}{2},\frac{1}{2},\frac{1}{4}$ 互相配套，并与费米子两点函数的归一化方式绑定：若要核对 (2.2)、(2.3) 的 OPE 系数，必须同时固定 $\psi$ 的归一化，单独改动其中任一系数都会破坏 $J$ 的 $U(1)$ 荷约定。此外，相互作用项只通过 $R_{i\bar{k}j\bar{l}}$ 出现，说明该 sigma 模型的相互作用完全由靶空间曲率决定。
\end{yizhu}
